% Generated from locale/tr/content/history/biographies/alfred-tarski.tex; do not hand-edit.


\olfileid[tr]{his}{bio}{tar}

\olsection{Alfred Tarski}

\olphoto{tarski-alfred}{Alfred Tarski}

Alfred Tarski, 14 Ocak 1901'de Polonya'nın Varşova kentinde (o sırada Rus
İmparatorluğu'nun parçasıydı) doğdu. ``Napolyonvari'' diye betimlenen Tarski
gürültücü, konuşkan ve yoğun biriydi. Enerjisi derslerine de sık sık yansırdı:
bir keresinde ders sırasında sigarasını atarken çöp sepetini tutuşturmuş ve o
binada bir daha ders vermesi yasaklanmıştı.

Tarski küçük yaştan beri bilgiye susamıştı. İlerleyen yıllarda öğrencilerine,
mantık okumasının nedeninin B aldığı tek dersin bu olması olduğunu anlatsa da
lise kayıtları mantık dâhil bütün derslerden A aldığını gösterir. 1918--1924
arasında Varşova Üniversitesi'nde öğrenim gördü. Tarski önce biyoloji okumayı
düşünüyordu; ancak üniversite Varşova Mantık ve Felsefe Okulu'nun merkezi
olduğundan matematik, felsefe ve mantığa ilgi duymaya başladı. Doktorasını
1924'te Stanis\l{}aw Le\'{s}niewski'nin danışmanlığında aldı.

Tarski, 1939'da Amerika Birleşik Devletleri'ne göç etmeden önce en önemli
çalışmalarından bazılarını Varşova'da ortaöğretim öğretmeni olarak çalışırken
tamamladı. Mantıksal sonuç ve mantıksal doğruluk üzerine eserleri bu dönemde
yazıldı. Tarski 1939'da bir konferans turnesi için Amerika Birleşik
Devletleri'nde bulunuyordu. Ziyareti sırasında Almanya Polonya'yı işgal etti;
Yahudi kökeni nedeniyle Tarski ülkesine dönemedi. Eşi ve çocukları savaşın
sonuna kadar Polonya'da kaldılar; daha sonra onlar da Amerika Birleşik
Devletleri'ne göç edebildi. Tarski Harvard'da, City College of New York'ta,
Princeton'daki İleri Araştırmalar Enstitüsü'nde ve son olarak Berkeley'deki
California Üniversitesi'nde ders verdi. Berkeley'de disiplinlerarası Mantık
ve Bilim Metodolojisi programını kurdu. Tarski 26 Ekim 1983'te 82 yaşında
öldü.

\begin{reading} 
Tarski'nin yaşamı hakkında daha fazla bilgi için \emph{Alfred Tarski: Life
  and Logic} adlı biyografiye bakın \citep{Feferman2004}. Tarski'nin mantıksal
sonuç ve doğruluk üzerine ufuk açıcı eserleri \citep{Tarski1983}'te İngilizce
olarak mevcuttur. Tarski'nin bütün özgün eserleri dört ciltlik bir dizide
toplanmıştır \citep{Tarski1981}.
\end{reading}

